\documentclass[11pt,a4paper]{article}

\usepackage[utf8]{inputenc}
\usepackage[T2A]{fontenc}
\usepackage[ukrainian,english]{babel}
\usepackage[top=22mm,bottom=22mm,left=22mm,right=22mm]{geometry}
\usepackage{hyperref}
\usepackage{fancyhdr}
\usepackage{parskip}
\usepackage{tcolorbox}
\usepackage{enumitem}
\usepackage{listings}
\usepackage{xcolor}

\tcbuselibrary{skins,breakable}

\hypersetup{
    pdftitle={Workshop 03 --- Slash Commands: Handout},
    pdfauthor={Vadym Bondarenko},
    colorlinks=true,
    linkcolor=blue,
    urlcolor=cyan,
}

\pagestyle{fancy}
\fancyhf{}
\fancyhead[L]{\small\textit{Workshop 03 --- Slash Commands}}
\fancyhead[R]{\small\thepage}
\fancyfoot[C]{\small\textit{vcdev.me \textperiodcentered{} 2026}}
\renewcommand{\headrulewidth}{0.4pt}

\lstdefinelanguage{yaml}{
    keywords={true,false,null},
    sensitive=false,
    comment=[l]{\#},
    morestring=[b]',
    morestring=[b]"
}

\lstset{
    basicstyle=\ttfamily\small,
    breaklines=true,
    frame=single,
    backgroundcolor=\color{gray!10},
    rulecolor=\color{gray!40},
    xleftmargin=8pt,
    xrightmargin=8pt,
    aboveskip=6pt,
    belowskip=6pt,
}

\title{%
  {\Huge\bfseries Workshop 03: Slash Commands}\\[0.6em]
  {\large Глибше у `/`-команди --- handout}
}
\author{Vadym Bondarenko \\ \small\href{mailto:vadym.bondarenko@vcdev.me}{vadym.bondarenko@vcdev.me}}
\date{2026}

\begin{document}

\maketitle

\begin{abstract}
Пост-воркшоп референс для самостійного перечитування. Не повторює слайди --- сухе зведення команд, frontmatter-полів, чек-листів і decision-trees. Усі факти посилаються на офіційні docs (\href{https://code.claude.com/docs/en/skills}{code.claude.com}).
\end{abstract}

\tableofcontents
\newpage

\section{The big merge: commands = skills}

З останнього оновлення Claude Code custom commands були злиті у формат skills.

\begin{tabular}{ll}
\hline
\textbf{Що} & \textbf{Шлях} \\
\hline
Legacy command (працює) & \texttt{.claude/commands/deploy.md} \\
New skill-form command   & \texttt{.claude/skills/deploy/SKILL.md} \\
\hline
\end{tabular}

Обидва створюють \texttt{/deploy}. \textbf{Skill виграє} при однаковому імені. Frontmatter --- спільний schema. Нові команди пишуться у skill-формі, бо отримують: supporting files, \texttt{disable-model-invocation}, \texttt{paths}, \texttt{context: fork}.

\section{Команда vs auto-skill}

Одна різниця у frontmatter:

\begin{lstlisting}[language=yaml]
---
disable-model-invocation: true   # це робить її командою
---
\end{lstlisting}

\begin{tabular}{lll}
\hline
\textbf{Frontmatter} & \textbf{Тригер} & \textbf{Use case} \\
\hline
default                            & \texttt{/name} АБО Claude          & Безпечні задачі \\
\texttt{disable-model-invocation: true} & \texttt{/name} тільки         & Side effects (deploy, commit) \\
\texttt{user-invocable: false}     & Claude тільки                      & Background знання \\
\hline
\end{tabular}

\section{Frontmatter cheatsheet}

\subsection{Командний мінімум}

\begin{lstlisting}[language=yaml]
---
name: git-summary                              # /git-summary
description: Summarize recent git commits      # текст у /-меню
argument-hint: [--since=<duration>]            # підказка автокомпліту
disable-model-invocation: true                 # робить це командою
allowed-tools: Bash(git log *)                 # pre-approval
---
\end{lstlisting}

\subsection{Усі релевантні поля}

\begin{tabular}{lll}
\hline
\textbf{Поле} & \textbf{Тип} & \textbf{Призначення} \\
\hline
\texttt{name} & string & \texttt{/name}; kebab-case, $\le$64 \\
\texttt{description} & string & UX у меню; $\le$1536 з \texttt{when\_to\_use} \\
\texttt{argument-hint} & string & автокомпліт-підказка \\
\texttt{arguments} & list & іменовані позиційні (\texttt{\$name}) \\
\texttt{disable-model-invocation} & bool & команда-only \\
\texttt{user-invocable} & bool & Claude-only (false = ховає) \\
\texttt{allowed-tools} & string/list & pre-approval \\
\texttt{paths} & glob list & обмеження автотригера \\
\texttt{shell} & string & \texttt{bash} (default) / \texttt{powershell} \\
\texttt{context: fork} & string & subagent-ізоляція \\
\texttt{model}, \texttt{effort} & string & override на час команди \\
\hline
\end{tabular}

\section{\$ARGUMENTS: чотири способи}

\begin{lstlisting}[language=yaml]
---
name: migrate-component
arguments: [name, from, to]
---
Migrate $name from $from to $to.
\end{lstlisting}

\textit{/migrate-component SearchBar React Vue}:

\begin{tabular}{ll}
\hline
\textbf{Синтаксис} & \textbf{Значення} \\
\hline
\texttt{\$ARGUMENTS} & \texttt{SearchBar React Vue} (все одним рядком) \\
\texttt{\$ARGUMENTS[0]} & \texttt{SearchBar} \\
\texttt{\$0}, \texttt{\$1}, \texttt{\$2} & \texttt{SearchBar}, \texttt{React}, \texttt{Vue} \\
\texttt{\$name}, \texttt{\$from}, \texttt{\$to} & те саме за іменами \\
\hline
\end{tabular}

\textbf{Лапки = один токен:} \texttt{/cmd "hello world"} $\to$ \texttt{\$0 = hello world}.

\textbf{Якщо плейсхолдера немає} --- Claude Code сам додає \texttt{ARGUMENTS: <text>} у кінець SKILL.md.

\section{Shell-ін'єкція через !`cmd`}

Виконується \textbf{до} того, як Claude побачить промпт. Вивід підставляється у тіло.

\subsection{Inline}

\begin{lstlisting}
- Working dir: !`pwd`
- Branch: !`git branch --show-current`
- Latest: !`git log -1 --oneline`
\end{lstlisting}

\subsection{Multi-line block}

\begin{lstlisting}
```!
node --version
git status --short
```
\end{lstlisting}

\subsection{Економія}

Без ін'єкції: 4--5 turn-ів (Claude $\to$ Bash $\to$ Claude $\to$ \dots). З ін'єкцією: 1 turn.

\subsection{Kill-switch}

\begin{lstlisting}
{ "disableSkillShellExecution": true }
\end{lstlisting}

Замінює кожен \texttt{!`cmd`} на \texttt{[shell command execution disabled by policy]}. Не зачіпає bundled і managed skills.

\section{allowed-tools: pre-approval, не sandbox}

\textbf{Що робить:} дозволяє listed tools без per-use approval prompt-у поки команда активна.

\textbf{Що НЕ робить:}
\begin{itemize}
\item НЕ блокує інші tool-и --- вони далі викликаються через approval prompt
\item НЕ ізолює від файлової системи
\item НЕ замінює permission settings
\end{itemize}

\textbf{Цитата docs:} \emph{``It does not restrict which tools are available: every tool remains callable, and your permission settings still govern tools that are not listed.''}

\subsection{Pattern syntax}

\begin{lstlisting}
allowed-tools: Bash(git add *) Bash(git commit *) Read
\end{lstlisting}

\begin{itemize}
\item \texttt{ToolName} --- весь tool
\item \texttt{ToolName(matcher)} --- glob-патерн
\item Space-separated string АБО YAML list
\end{itemize}

\subsection{Справжній sandbox через settings.json}

\begin{lstlisting}
{
  "permissions": {
    "allow": [
      "Bash(git status *)",
      "Bash(git log *)"
    ],
    "deny": [
      "Bash(rm *)",
      "Bash(curl *)",
      "Bash(git push *)"
    ]
  }
}
\end{lstlisting}

\textbf{Combo} \texttt{allowed-tools} (UX) + \texttt{deny rules} (security) = щільний контроль.

\section{Plugin packaging}

\subsection{Структура}

\begin{lstlisting}
my-git-commands/
├── .claude-plugin/
│   └── plugin.json     # тільки маніфест тут
└── skills/
    ├── git-summary/SKILL.md
    └── env-info/SKILL.md
\end{lstlisting}

\textbf{Pitfall:} НЕ кладіть \texttt{skills/}, \texttt{commands/}, \texttt{agents/}, \texttt{hooks/} \textbf{всередину} \texttt{.claude-plugin/}.

\subsection{Мінімальний plugin.json}

\begin{lstlisting}
{
  "name": "my-git-commands",
  "description": "Git workflow commands",
  "version": "0.1.0",
  "author": { "name": "You" }
}
\end{lstlisting}

\subsection{Namespacing}

\begin{tabular}{ll}
\hline
\textbf{Що} & \textbf{Інвокація} \\
\hline
Personal/Project skill & \texttt{/git-summary} \\
Plugin skill           & \texttt{/my-git-commands:git-summary} \\
\hline
\end{tabular}

Префікс \texttt{my-git-commands} = поле \texttt{name} у \texttt{plugin.json}.

\subsection{Локальне тестування}

\begin{lstlisting}
claude --plugin-dir ./my-git-commands
\end{lstlisting}

\begin{itemize}
\item Завантажує плагін без installation
\item Перебиває одноіменний marketplace плагін на сесію
\item Multiple flags: \texttt{--plugin-dir ./p1 --plugin-dir ./p2}
\item \texttt{/reload-plugins} підхоплює зміни SKILL.md
\end{itemize}

\subsection{Версіонування}

\begin{itemize}
\item Set explicit \texttt{version} $\to$ юзери оновлюються лише при bump-і
\item Skip $\to$ commit SHA = версія, кожен коміт оновлює юзерів
\end{itemize}

\section{Debug: 5-крокова перевірка}

\begin{enumerate}
\item \textbf{Шлях точний?}
  \begin{lstlisting}
ls ~/.claude/skills/<name>/SKILL.md
ls .claude/skills/<name>/SKILL.md
  \end{lstlisting}

\item \textbf{Frontmatter валідний YAML?} \texttt{---} зверху і знизу, відступи пробілами

\item \textbf{\texttt{name} валідне?} kebab-case, цифри, дефіси, $\le$64 символів

\item \textbf{Створював нову теку під час сесії?} Потрібен рестарт Claude Code

\item \textbf{Плагін?} \texttt{--plugin-dir} був переданий? \texttt{/reload-plugins} після edit?
\end{enumerate}

\subsection{Часті проблеми}

\begin{tabular}{ll}
\hline
\textbf{Симптом} & \textbf{Перевірка} \\
\hline
Команди нема у \texttt{/}-меню & Шлях, frontmatter, рестарт \\
\texttt{\$ARGUMENTS} лишився як literal & Синтаксис підстановки, YAML парсинг \\
\texttt{!`cmd`} поверне placeholder & \texttt{disableSkillShellExecution} у settings \\
Description обрізається & \texttt{SLASH\_COMMAND\_TOOL\_CHAR\_BUDGET} \\
\hline
\end{tabular}

\section{Production-чек-ліст}

Перед публікацією:

\begin{itemize}[leftmargin=2em]
\item[$\square$] \texttt{name} kebab-case, унікальний
\item[$\square$] \texttt{description} однорядкова, корисна у \texttt{/}-меню
\item[$\square$] \texttt{argument-hint} для команд з аргументами
\item[$\square$] \texttt{disable-model-invocation: true} для всього з side effects
\item[$\square$] \texttt{allowed-tools} лише потрібне
\item[$\square$] \texttt{!`cmd`} лише там, де економія турнів варта
\item[$\square$] Deny-rules у settings.json для справжньої ізоляції
\item[$\square$] Тестовано на чистій сесії
\item[$\square$] Plugin: \texttt{version} у \texttt{plugin.json} для semver
\item[$\square$] Plugin: \texttt{README.md} як встановити, що робить
\end{itemize}

\section{Команди для вправ}

\subsection{Вправа 1 --- встановити команду}

\begin{lstlisting}
mkdir -p ~/.claude/skills/git-summary
cp starter/SKILL.md ~/.claude/skills/git-summary/SKILL.md
# рестарт Claude Code
/git-summary --since=yesterday
\end{lstlisting}

\subsection{Вправа 2 --- shell-ін'єкція}

\begin{lstlisting}
mkdir -p ~/.claude/skills/env-info
# у SKILL.md:
- Working dir: !`pwd`
- Branch: !`git branch --show-current`
\end{lstlisting}

\subsection{Вправа 3 --- allowed-tools}

\begin{lstlisting}
allowed-tools: Bash(git status *) Bash(git log *) Bash(git branch *)
# Бонус: deny-rules у ~/.claude/settings.json
\end{lstlisting}

\subsection{Вправа 4 --- спакувати у плагін}

\begin{lstlisting}
mkdir -p my-git-commands/.claude-plugin
cp -r starter/skills my-git-commands/skills
# написати plugin.json
claude --plugin-dir ./my-git-commands
/my-git-commands:git-summary --since=yesterday
\end{lstlisting}

\section{Resources}

\begin{itemize}
\item \href{https://code.claude.com/docs/en/skills}{code.claude.com/docs/en/skills} --- Skills + commands merged
\item \href{https://code.claude.com/docs/en/plugins}{code.claude.com/docs/en/plugins} --- Plugin manifest, namespacing
\item \href{https://code.claude.com/docs/en/permissions}{code.claude.com/docs/en/permissions} --- Allow/deny rules
\item \href{https://code.claude.com/docs/en/commands}{code.claude.com/docs/en/commands} --- Built-in commands reference
\item Цей репо: \texttt{workshops/03-slash-commands/exercises/} і \texttt{solutions/}
\end{itemize}

\end{document}
